\section{Thử thách và cách giải quyết}
\label{sec:thu-thach}

Trong quá trình cài đặt từ những dòng socket đầu tiên cho tới khi ghép toàn bộ hệ thống, nhóm đã vấp phải khá
nhiều bug mà lúc chỉ vẽ thiết kế trên giấy hoặc chạy thử vài gói trên loopback không hề thấy. Đa số các lỗi khó
đều chỉ bộc lộ khi nhóm đưa vào chạy tải nặng trên bộ giả lập mạng có rớt gói và độ trễ cao. Chương này ghi lại
những vấn đề đáng nhớ nhất ở cả hai tầng cùng cách nhóm tìm nguyên nhân và xử lý.

\subsection{Tầng giao thức và transport}
\begin{enumerate}[label=\textbf{\arabic*.}, ref=\arabic*, leftmargin=*, itemsep=4pt]
  \item \label{tt:cubic-cham}\textbf{Tốc độ truyền với CUBIC bị tụt tới 280 lần khi bật loss 10\%.} Khi in log giá
        trị \texttt{cwnd} ra kiểm tra, nhóm thấy cửa sổ gần như luôn bị ghim ở mức sàn 2 gói. Với chỉ 2 gói được
        phép bay trên mạng cùng lúc, bên gửi không bao giờ nhận được đủ 3 gói đi sau vượt mặt để kích hoạt truyền
        lại nhanh, dẫn tới mọi lần rớt gói đều phải chờ hết bộ hẹn giờ RTO. Tệ hơn nữa, lúc đầu nhóm đặt
        RTO tối thiểu bằng 200 ms theo thói quen của TCP trong khi RTT thực trên loopback chỉ khoảng 0,4 ms.
        \emph{Cách giải quyết:} nhóm bổ sung thêm điều kiện phát hiện mất gói theo ngưỡng thời gian (sau
        $\tfrac98\,\text{RTT}$ kể từ lúc gửi nếu đã có gói đi sau được ACK), đồng thời chuyển sang công thức PTO
        của QUIC \cite{rfc9002} và bỏ mức sàn 200 ms do bên nhận của nhóm luôn trả ACK ngay lập tức (\cref{subsec:mat-goi}).
  \item \label{tt:jitter}\textbf{Có tới 89\% số gói bị báo mất nhầm dù cấu hình \texttt{loss = 0\%}.} Khi mới bật
        thử tham số \texttt{Jitter}, nhóm thấy bộ đếm gói truyền lại tăng vọt dù không hề bật rớt gói. Kiểm tra
        kỹ thì lỗi nằm ở chính bộ giả lập \fname{LossyConn} chứ không phải ở giao thức: bản đầu tiên của nhóm dùng
        \fname{time.AfterFunc} hẹn giờ gửi riêng rẽ cho từng gói với độ trễ ngẫu nhiên, vô tình biến đường truyền
        thành đảo thứ tự gói hàng loạt, làm quy tắc ``3 gói đi sau tới trước'' liên tục báo mất gói giả. Trong thực
        tế, các gói đi qua cùng một hàng đợi router vẫn giữ thứ tự FIFO. \emph{Cách giải quyết:} nhóm sửa lại bộ
        giả lập jitter đẩy gói qua một hàng đợi FIFO chung, sau đó tỉ lệ báo mất nhầm giảm từ 89\% xuống chỉ còn
        0,1--0,2\%. Đây là bài học lớn cho nhóm về việc phải viết unit test kiểm tra cả công cụ test trước khi tin
        vào số liệu đo.
  \item \label{tt:ping-dau}\textbf{Client vừa kết nối xong bị ``treo'' 2--4 giây đầu nếu mạng có rớt gói.} Do
        server chỉ học được địa chỉ IP:port UDP của client khi nhận được gói UDP đầu tiên từ client gửi lên, nếu
        chẳng may gói PING mở màn này bị rớt trên đường truyền thì server hoàn toàn không biết gửi dữ liệu cho
        client về địa chỉ nào, phải đợi tới chu kỳ keep-alive 2 giây sau mới có gói PING tiếp theo. \emph{Cách
        giải quyết:} khi mới mở kết nối và chưa nhận được gói nào từ server, client chủ động gửi lại gói PING đầu
        tiên theo chu kỳ ngắn 200, 400, 800 ms\dots{} (tương tự cơ chế truyền lại gói Initial trong QUIC \cite{rfc9002}). Nhờ vậy
        thời gian hoàn tất truyền một file 1 KB trên mạng loss 10\% giảm từ 2--4 giây xuống dưới 0,6 giây.
  \item \label{tt:cpu}\textbf{Tiến trình server ngốn 43\% CPU dù 20 client đang đứng im không gửi gì.} Khi chạy
        thử tải 20 client ở trạng thái rảnh, nhóm bật \texttt{pprof} lên xem thì thấy phần lớn thời gian CPU nằm ở
        \texttt{futex} và \texttt{findRunnable} của runtime Go. Nguyên nhân là mỗi kết nối đều giữ một
        \fname{time.Ticker} gõ đều đặn mỗi 10 ms bất kể có dữ liệu hay không, khiến bộ lập lịch goroutine bị đánh
        thức hàng nghìn lần mỗi giây. \emph{Cách giải quyết:} nhóm đổi sang cơ chế hai tốc độ: \texttt{ticker} chỉ
        chạy ở chu kỳ 10 ms khi bảng \texttt{unacked} đang có gói chờ ACK, còn khi đã ACK hết thì tự giãn ra
        100 ms (vì các tác vụ lúc rảnh như keep-alive 2 s hay idle timeout 10 s không cần độ phân giải tới 10 ms).
        Sau khi sửa, mức chiếm dụng CPU lúc rảnh giảm xuống còn 3,1\% mà không ảnh hưởng tới thông lượng.
  \item \textbf{Đồ thị \texttt{cwnd} của CUBIC lúc đầu có dạng răng cưa tuyến tính giống hệt Reno.} Khi xuất dữ
        liệu ra vẽ đồ thị \cref{fig:cwnd} lần đầu, nhóm phát hiện các đoạn tăng cửa sổ hoàn toàn thẳng chứ không
        có đoạn đi ngang bậc ba quanh $W_\text{max}$. Lý do là trên mạng có RTT ngắn, nhánh tăng cộng tính kiểu
        Reno đã lấn át hoàn toàn hàm bậc ba do cách nhóm viết điều kiện rẽ nhánh ban đầu chưa chuẩn. \emph{Cách
        giải quyết:} nhóm cài đặt lại hàm \fname{OnAck} bám sát RFC~9438 \cite{rfc9438}, cập nhật song song hai biến
        $W_\text{cubic}$ và $W_\text{est}$ rồi lấy $\max(W_\text{cubic}, W_\text{est})$, sau đó vẽ lại đồ thị
        phóng to (\cref{fig:cwnd-zoom}) để xác nhận đường cong chữ S bậc ba.
  \item \textbf{Tin nhắn chat hiển thị sai thứ tự dù mạng không rớt gói.} Do nhóm tách mỗi tin nhắn ra một stream
        độc lập để tránh HoL, khi người dùng gửi liên tiếp nhiều tin ngắn thì tùy vào cách ghép gói mà stream sau
        có thể hoàn tất và đẩy lên channel \texttt{Recv()} trước stream trước. \emph{Cách giải quyết:} nhóm bổ sung
        trường số thứ tự \texttt{Seq} ở bản tin tầng ứng dụng và đặt bộ đệm \fname{SeqBuffer} để sắp lại đúng thứ
        tự gửi, nhờ đó vừa giữ được tính chất không HoL giữa file với chat ở tầng transport, vừa đảm bảo thứ tự
        hội thoại cho người dùng (\cref{subsec:app}).
  \item \label{tt:nonce}\textbf{Trùng lặp nonce AES-GCM giữa chiều gửi của client và chiều gửi của server.} Ở
        phiên bản đầu, cả hai phía cùng dùng khóa $k_\text{udp}$, cùng khởi tạo \texttt{packet\_num = 0} và ghép
        nonce chỉ từ \texttt{(conn\_id, packet\_num)}, dẫn tới gói số $n$ từ client gửi lên và gói số $n$ từ
        server gửi xuống dùng chung một cặp (khóa, nonce). \emph{Cách giải quyết:} nhóm đưa thêm 1 byte định danh
        chiều truyền (\texttt{0} cho C$\to$S, \texttt{1} cho S$\to$C) vào byte thứ 9 của nonce và viết unit test
        kiểm tra chéo (\cref{sec:bao-mat}).
  \item \textbf{Nhập nhằng mẫu đo RTT khi gói bị truyền lại.} Trong bài toán kinh điển của TCP, khi một segment
        phải truyền lại thì lúc nhận ACK về sẽ không phân biệt được ACK đó trả lời cho lần gửi gốc hay lần gửi
        lại. \emph{Cách giải quyết:} nhóm áp dụng triệt để nguyên tắc của QUIC \cite{rfc9002} là không bao giờ tái sử dụng
        \texttt{packet\_num}; mỗi gói truyền lại đều mang số \texttt{packet\_num} mới và ghi nhận mốc thời gian
        gửi mới, nên phép đo RTT luôn chính xác mà không cần tới thuật toán Karn \cite{rfc6298}.
  \item \textbf{Ghi vào stream của kết nối đã đóng nhưng hàm \texttt{Write} vẫn báo thành công.} Trong câu lệnh
        \fname{select} của Go, nếu nhánh đẩy dữ liệu vào channel (vẫn còn chỗ trống trong buffer) và nhánh
        \texttt{<-c.done} cùng thỏa mãn tại một thời điểm, runtime Go sẽ chọn ngẫu nhiên một trong hai nhánh.
        \emph{Cách giải quyết:} nhóm thêm một bước kiểm tra nhanh \texttt{select \{ case <-c.done: return err;
        default: \}} ngay ở đầu hàm trước khi thực hiện \fname{select} ghi vào channel.
  \item \textbf{Vòng lặp ACK vô tận và phân mảnh gói IP.} Ở bản thử nghiệm đầu tiên, hễ nhận được gói UDP nào là
        bên nhận lại đánh dấu cần gửi ACK, khiến gói chỉ chứa ACK cũng kích hoạt một gói ACK trả lời ngược lại
        mãi không dứt. \emph{Cách giải quyết:} nhóm quy định chỉ bật cờ \texttt{needAck} khi trong gói nhận được có
        ít nhất một frame STREAM hoặc PING, đồng thời khóa trần kích thước gói ở mức 1200 byte để tránh phân mảnh
        IP.
  \item \textbf{Data race giữa luồng đọc socket và luồng của ứng dụng.} Khi chạy \fname{go test -race} ở những bản
        đầu, trình kiểm tra báo xung đột bộ nhớ khi goroutine ứng dụng đọc thông số kết nối cùng lúc goroutine nhận
        gói đang cập nhật state. \emph{Cách giải quyết:} nhóm chuyển hẳn sang kiến trúc một goroutine
        \fname{run()} duy nhất sở hữu toàn bộ state của \fname{Conn} và chỉ trao đổi với bên ngoài qua channel
        (\cref{subsec:concurrency}).
\end{enumerate}

\subsection{Tầng ứng dụng}
\begin{enumerate}[label=\textbf{\arabic*.}, leftmargin=*, itemsep=4pt]
  \item \textbf{Tin nhắn đến người nhận bị đảo thứ tự do gọi \texttt{go Send()} tràn lan.} Ban đầu, để vòng lặp
        đọc của server không bị chặn khi broadcast tin nhắn cho nhiều người, nhóm tạo một goroutine mới cho mỗi
        lần gửi (\texttt{go session.Send(msg)}). Vì bộ lập lịch của Go không đảm bảo goroutine nào tạo trước sẽ
        chạy trước, khi gửi 20 tin liên tiếp thì phía nhận thường xuyên thấy tin số 2 hiện lên trước tin số 1 ngay
        cả khi mạng không mất gói nào. \emph{Cách giải quyết:} mỗi session được cấp một hàng đợi \texttt{sendCh}
        riêng và chỉ có duy nhất một goroutine \fname{writeLoop} lấy tin ra gọi \texttt{SendMessage} tuần tự, kèm
        test tự động gửi 20 tin liên tiếp để kiểm tra thứ tự.
  \item \textbf{Một người nhận đọc chậm hoặc bị treo làm kẹt toàn bộ việc định tuyến trên server.} Nếu server vừa
        giữ khóa \fname{Registry} vừa gọi hàm gửi mạng, hoặc đẩy vào channel \texttt{sendCh} theo kiểu chặn
        (blocking), chỉ cần một client bị đơ (làm cửa sổ nhận về 0 và \texttt{sendCh} đầy) là goroutine định tuyến
        của người gửi sẽ bị treo theo, kéo sập cả phòng chat. Ở phiên bản đầu, nhóm thử sửa bằng cách bỏ qua tin mới
        khi hàng đợi đầy, nhưng cách này làm mất tin nhắn một cách lặng lẽ và để lại khoảng trống \texttt{Seq} khiến
        \fname{SeqBuffer} ở client đứng chờ. \emph{Cách giải quyết:} server chỉ giữ khóa \fname{Registry} để chép
        nhanh danh sách con trỏ session ra ngoài (\fname{Snapshot}) rồi nhả khóa ngay trước khi gửi; thao tác đẩy
        vào \texttt{sendCh} (4096 tin) dùng \texttt{select} có nhánh \texttt{default}, và nếu hàng đợi của client
        nào bị đầy tràn thì server chủ động ngắt kết nối của riêng client quá chậm đó để client tự kết nối lại
        (\cref{subsec:app}).
  \item \textbf{Lọt hai phiên trùng username khi hai client đăng nhập cùng lúc.} Ban đầu nhóm viết hai bước tách
        rời: gọi \fname{GetByName} để kiểm tra xem tên đó đã online chưa, nếu chưa thì mới gọi hàm thêm vào
        \fname{Registry}. Khi test thử trường hợp hai client dùng cùng một tài khoản gửi bản tin \texttt{AUTH}
        đồng thời, cả hai goroutine đều đi qua bước kiểm tra trước khi bên kia kịp ghi vào map. \emph{Cách giải
        quyết:} nhóm gộp bước kiểm tra và bước đăng ký vào chung một hàm nguyên tử \fname{AddPending} dưới cùng
        một lần giữ \texttt{Lock()}, trả về lỗi \texttt{ErrAlreadyOnline} ngay nếu tên đó đã tồn tại.
  \item \textbf{Xóa nhầm phiên mới khi người dùng rớt mạng rồi kết nối lại nhanh.} Trước đây hàm dọn dẹp
        \fname{dropSession} chỉ xóa theo chuỗi \texttt{username}. Khi một client bị rớt mạng (server chưa hết 10 s
        idle timeout) và người dùng khởi động lại client để đăng nhập phiên mới, một lát sau kết nối cũ hết
        timeout mới gọi hàm xóa theo tên và xóa mất luôn phiên mới vừa tạo, biến client mới thành trạng thái
        ``bóng ma'' (vẫn nhận được tin nhưng biến mất khỏi danh sách online). \emph{Cách giải quyết:} hàm
        \fname{Registry.Remove(s *Session)} so sánh trực tiếp con trỏ \fname{Session} và chỉ xóa khỏi map nếu con
        trỏ đang lưu vẫn đúng là phiên \texttt{s} hiện tại của tên đó, đồng thời thiết kế idempotent để gọi từ nhiều
        nơi (khi vòng đọc thoát, khi \texttt{Done()} đóng, hoặc khi gặp lỗi ghi) mà không gây tác dụng phụ.
  \item \textbf{Phiên pending bị treo vĩnh viễn nếu client mất mạng ngay sau \texttt{AUTH\_OK}.} Nếu một
        client hoàn tất handshake TCP nhưng chưa kịp gửi gói UDP nào đã bị mất kết nối, phiên \textsc{Pending} sẽ
        giữ tên tài khoản đó ở trạng thái đang online mãi mãi khiến chính người đó không thể đăng nhập lại. Lúc đầu
        nhóm định viết thêm một goroutine quét dọn timer riêng cho các phiên pending; \emph{cách giải quyết} gọn
        hơn là gọi \fname{Mux.Accept} tạo sẵn đối tượng \fname{transport.Conn} ngay từ lúc trả \texttt{AUTH\_OK}:
        khi đó nếu sau 10 giây không có gói UDP nào tới thì chính cơ chế idle timeout có sẵn của tầng transport sẽ
        tự đóng kết nối và kích hoạt dọn dẹp session mà không tốn thêm bộ hẹn giờ mới.
  \item \textbf{Chuyển tiếp file kích thước lớn qua server mà không làm tràn RAM.} Với kịch bản 20 client gửi file
        200 MB cho nhau, nếu server gom toàn bộ file vào bộ nhớ rồi mới gửi tiếp cho người nhận thì server sẽ nhanh
        chóng cạn RAM. \emph{Cách giải quyết:} nhóm thiết kế cơ chế chuyển tiếp kiểu pipeline: ngay khi
        \fname{FileHeaderParser} gom và tách xong phần JSON metadata ở đầu stream (xử lý được cả trường hợp JSON bị
        cắt ngang giữa hai mảnh \fname{Delivery} đầu tiên), server mở ngay một stream hướng tới người nhận và đẩy
        từng mảnh dữ liệu qua một channel đệm 256 phần tử cho goroutine \fname{relayFile}. Khi người nhận đọc chậm
        làm channel này đầy, vòng đọc phía người gửi sẽ bị chặn lại và cơ chế điều khiển luồng của tầng transport sẽ
        tự động điều tiết bên gửi chậm theo.
  \item \textbf{Tự động kết nối lại và đóng kết nối gọn khi bấm Ctrl+C.} Khi channel \texttt{Done()} đóng mà không
        phải do người dùng chủ động thoát, client chuyển sang trạng thái \textsc{Reconnecting} với thời gian chờ
        backoff tăng dần 1, 2, 4\dots{} tới tối đa 30 s (nếu trong 10 s đầu server chưa kịp dọn phiên cũ và trả về
        \texttt{AUTH\_FAIL} báo đang online thì các lần thử kế tiếp sẽ tự vào lại được). Ngoài ra, ở các bản đầu khi
        bấm Ctrl+C để tắt, tiến trình bị dừng ngay mà không kịp gửi frame \texttt{CLOSE} nên những người khác phải
        chờ hết 10 s mới thấy danh sách online cập nhật. \emph{Cách giải quyết:} nhóm bắt tín hiệu
        \texttt{SIGINT}/\texttt{SIGTERM} ở cả client lẫn server (phía server đóng \fname{Mux} để phát \texttt{CLOSE}
        cho toàn bộ client đang kết nối) trước khi thoát, rút ngắn thời gian cập nhật trạng thái xuống chỉ còn
        0,05--0,2 giây (\cref{tab:he-thong}).
  \item \textbf{Tự động tìm server trong mạng LAN.} Ở phiên bản đầu, gói UDP broadcast ghi cứng cổng 9000 kèm địa
        chỉ IP do server tự đoán (dẫn tới lấy sai IP khi máy chạy server có nhiều card mạng ảo). \emph{Cách giải
        quyết:} server phát đúng số cổng TCP thực tế đang lắng nghe, còn phía client lấy IP của server trực tiếp từ
        địa chỉ nguồn của gói UDP broadcast nhận được. Bên cạnh đó, để nhiều tiến trình client trên cùng một máy có
        thể cùng lắng nghe cổng 9002 mà không bị lỗi trùng bind ở tiến trình thứ hai, nhóm bật tùy chọn
        \texttt{SO\_REUSEADDR}/\texttt{SO\_REUSEPORT} trên socket (tách file cài đặt riêng cho Windows và Unix thông
        qua build tag).
  \item \textbf{Cờ giả lập mạng được truyền vào nhưng không có tác dụng.} Ở một phiên bản thử nghiệm, cờ dòng lệnh
        \texttt{-jitter} đã được hàm \texttt{main} đọc vào nhưng lại quên gán xuống cấu hình của \fname{LossyConn};
        lỗi này chỉ lộ ra khi nhóm kiểm tra thấy RTT đo được vẫn giữ nguyên ở mức 0,1 ms. \emph{Cách giải quyết:}
        nhóm nối lại đầy đủ các tham số cấu hình và rút ra bài học tương tự thử thách~\ref{tt:jitter}: luôn phải
        kiểm tra bộ giả lập bằng số đo thực tế (RTT, tỉ lệ mất gói) trước khi tin vào kết quả đo.
  \item \textbf{Tự cài đặt HTTP server cho Web UI khi không được dùng \fname{net/http}.} Vì phải làm việc trực tiếp
        trên socket TCP, nhóm tự xử lý việc đọc request line và từng dòng header, dùng \texttt{io.LimitReader} theo
        đúng \texttt{Content-Length} cho cả request JSON nhỏ lẫn file upload hàng trăm MB (phần body của file được
        stream thẳng từ socket ra file tạm trên đĩa chứ không đọc dồn vào RAM), đồng thời trả header
        \texttt{Connection: close} để đơn giản hóa vòng đời kết nối HTTP. Trên trình duyệt, do server không gửi dội
        lại tin nhắn về cho chính người gửi nên phía JS tự hiển thị tin ngay khi bấm gửi; toàn bộ nội dung chat từ
        người khác gửi tới đều đi qua hàm escape ký tự HTML trước khi chèn vào DOM để chống lỗi XSS; ngoài ra khi
        chạy client ở chế độ nền (\texttt{< /dev/null}), dù luồng đọc \texttt{stdin} chạm \texttt{EOF} ngay lập tức
        nhưng tiến trình vẫn được giữ hoạt động chừng nào Web UI còn đang chạy.
\end{enumerate}
